% ================================================================
%  第 8 章 师与生 (原书 p.175-186, 后半: 8.5-8.8)
%  源文: _work/src1/task-ch08b.md
%  说明: 本章前半 (8.1-8.4) 见 part1-ch08a.tex
% ================================================================

% ---------------- p.175 ----------------
\section{(29) 敌父(1)}

1970年那个转折之后,在完全不同的环境里——一所外省大学——开始跟我做工作的那类学生,与从前的学生也很不相同。此外,在我这里做到国家博士论文水平的,只有两个人。其余人的工作,都在DEA或第三阶段博士论文的水平上。我还该算上一大批紧紧跟住某些“入门研究”课程的学生:这些课对他们来说,是提出往往出人意料的数学问题的机会,有时也是想象出别出心裁的方法去解决它们的机会。我遇到的最积极的参与,是在给一年级学生开的某些“选修课”上。相反,在那些已经在大学氛围里泡了几年的学生身上,某种清新、某种兴趣的能力、某种个人的眼光,已经或多或少褪了色。选修课的学生里,有几个显然具备成为出色数学家的素质。看当时的情形,我克制住自己,没有鼓励任何一个人走上那条路——尽管那条路本可以吸引他们,他们本可以在那里出类拔萃。

我跟那些为准备硕士文凭而听我这类“课程”的学生,关系大多没有延续到一年以上。每一次,我的印象都是:这些关系很快变得亲切而松弛,总的说来如此。只有一个被无孔不入的“怯场”缠住的学生是例外\fnoteP{9}{(23$''$)}{\emph{不敢玩}\par 这个学生曾和我一起做了一整年的DEA“实习工作”,而在他与我的工作关系里,直到最后都保持着“紧绷”状态。那是一种坦率友好的关系,贯穿着一种无可怀疑的相互好感。然而可就是有这份“怯场”;这份恐惧——它真正的起因肯定不是对我这个人的畏惧,尽管看上去像是如此。要不是这个学生自己跟我谈起,我也许根本不会察觉这件事;他谈,无疑是为了或多或少“解释”一下他在这一年里工作几乎完全卡住的原因。\par 正如在其他学生身上发生过的那样——他们和他一样,起初对某种几何实质抓得很牢——卡住恰恰出现在要“做零件活儿”的时候:也就是把成形的陈述白纸黑字写下来,或者仅仅是把握我提供的那些陈述的意义与含义;我提议把它们当作一种语言的基石、当作“游戏规则”接受下来。“学校式的”反应几乎总是把面对“做研究”局面的学生推向:把导师传下来的那些心照不宣的“游戏规则”当作既定的、又必须服从的东西接受下来,而这些规则恰恰不是要试着说明白的,更不是要弄懂的。这些心照不宣的规则所采取的具体形式,就是语义学或计算上的“配方”,不妨说,像预科班的题解书那样(或者任何一本普通教科书)。此外,学生还期待导师给出“证明……”这种形式的任务——那是他以往经验里遇见的唯一一种数学“思考”形式。(而且我不认为大多数职业数学家、以及其他科学家的禀性有本质不同——只不过“导师”换成了“共识”:共识规定当下的游戏规则,并把它当作不可更动的既定之物。这共识还规定哪些“问题”是要去解决的;在这些问题之间,每个人都觉得自己有随意挑选的自由,甚至允许自己在工作过程中修改它们,乃至发明另一些……)。我注意到,我面对有待探究的数学实质时那种完全不同的态度——因而也是面对学生时那种完全不同的态度——几乎必定会引发一种失措,而焦虑就是它的征兆之一。像一切焦虑一样,它倾向于给自己找一张面孔,投射到某个外部的“理由”上,不管那理由说得通说不通。焦虑最常见的一张面孔,恰恰就是恐惧。\par 这类困难在我教学生涯的第一个时期几乎没有出现,也许只有两例例外:那两次“师生”关系没超过几个星期;也许还有(我说不准)“忧伤的学生”那一例——他也许感到自己被“铆”在一个完全不启发他的题目上,而他本来完全有更换的自由。在那个(我也谈到过的)长期被某种怯场困扰的学生身上,原因显然在别处。他的工作丝毫没有卡住,恰恰相反,他对选定的主题极为自在,并在其上做出了一项有规模的基础性工作。那个时期我的学生大多是巴黎高师的旧生,他们与亨利·嘉当(Henri Cartan)的接触,已经让他们见识了另一种“接近数学”的方式。与之(可以说)相反的另一极,是我在蒙彼利埃大学任教的第二个时期:正是在一年级学生那里,我所说的焦虑最少干扰思考性的工作。在这些学生中的许多人身上,面对一种不同进路时的惊讶既不引起焦虑,也不引起封闭,相反带来的是敞开和兴致——终于可以做点有趣的事了!据我的观察,几年大学对一个学生创造性禀性的作用是彻底而毁灭性的。奇怪的是,在这方面,长年中学的作用似乎相对无害。原因也许是:\par 大学那几年正落在这样一个年纪:我们身上与生俱来的创造力,最迟必须通过个人的工作表达出来,否则就要永远沉船——至少就智力上的创造性工作而言是如此。肯定是出于一种健康的本能,我在做学生的那些年(也是在蒙彼利埃大学),几乎不踏进课堂,把几乎全部精力都投给了个人的数学思考。} (23'),
% ---------------- p.176 ----------------
\origpage{67} 其余则都是如此——那些本应正式在我指导下准备研究工作的学生,无论哪一个层次,也一样。与从前的学生相比,一个不同之处(还有许多别的!)是:我们的关系不那么局限于共同的数学工作。学生与我之间的交流,常常以不那么表面的方式牵涉到我们两个人本身\fnoteP{10}{(23$^{v}$)}{这种差别有一个格外醒目的征兆,出现在“外国人事件”之际——我曾有机会谈到过(第24节)。当时,许多与我完全陌生的人向我表达了同情,而我记不起1970年之前的哪个学生想到过以这种方式表示什么,更不用说在我投身的那场行动中向我提供任何帮助。相反,在我看来,第二个时期的学生或旧学生里,没有一个不曾向我表达同情与声援,还有几个积极参加了我在当地开展的运动。在这个狭小圈子之外,1945年法令那件事也在许多仅仅知道我名字的大学生中间激起了一定的波澜;在我被传唤的那一天,有不少人来到法院,以示声援。这后一个情形还提示:我所察觉到的“之前”与“1970年之后”学生态度上的差别,也许较少表达他们与我之间关系的差别,而更多表达的是一种心态的差别。显然,我“之前”的那些学生已经成了重要人物,而要让重要人物动心,是需要很多条件的……但我1970年离开IHÉS、投身战斗行动那一幕,似乎表明并不只是这一点。那是个他们中间还没有谁太像重要人物的时刻,然而我记不起他们中有谁对我所投身的活动表现出丝毫兴趣。我倒认为,那一定让他们所有人,无一例外地感到不自在。这仍然合乎“心态差别”的方向,但这不能只归因于社会地位的差别。} (23'')。因此,在我教学生涯的第二个时期,与某些学生的关系中的冲突成分表现得更清楚、更直接,甚至更激烈,这就不奇怪了。在第一个时期的旧学生里,有两个人后来表现出系统而毫不含糊的对抗态度(我曾顺带提到过),不过始终停留在未言说的层面,也许甚至停留在无意识的层面。在更长一些的第二个时期,有三个学生让我遭遇了对抗。其中两个,表现得相当尖锐。

这些学生中的一个,对抗是在一段最友好的关系里一夜之间冒出来的——那是在这位朋友不再是我的学生许多年之后。我猜测,冲突的原因与其说是我那不堪形容的举止和人格,不如说是一份压抑已久的不满:他的工作(曾经很出色)没有得到他本有权期待的那份接待。这就是“1970年之后”以我为老板的那份可疑特权的背面;他必定为此怨我,尽管即便在内心深处也不大肯向自己承认。

另一个学生,尖锐的对抗在工作一年半之后就已经出现,而此前的氛围看上去非常亲切。这是第一次也是唯一一次,学生与我之间的关系困难出现在他还是学生身份的时候。它使共同工作的继续成为不可能——而那份工作本来开局吉利,带着最好的兆头般的热情,而且应当说,题目是一个绝好的思考主题。我的感觉是:这位年轻研究者身上有一种暗中作祟的自信不足——不相信自己有能力做出好工作(这能力在我看来毫无疑问);而对抗以尖锐的音调爆发出来,是某种“向前逃窜”,为的是抢在一种他惧怕的失败之前动手,并把责任预先推到一位可憎的老板头上\fnoteP{11}{(23$'''$)}{\emph{两兄弟}\par 这个学生身上的对抗,一上来就采取了“阶级对抗”的形式:我是那个对他的数学前途拥有“生杀之权”的“老板”,可以凭自己的好恶决定……当然,事情本身只能证实这种看法,因为我很快就不再承担指导这个学生的责任(那已变得令人难受)。这把他置于一个微妙的处境——在眼下的世道里,找一个“老板”可不是那么显而易见的事,尤其是题目已经选定的时候。在另一个学生身上,合法期待落了空,对抗采取了类似的形式。我被感受为专横的“学阀”,不能容忍那些被他视为下属的人(学生或级别较低的同事)有任何反驳。\par 这样一种“阶级态度”,在我与第一个时期学生的关系里,从未有过一丝一毫的表现。显而易见的原因是:在1970年以前的形势下,学生通过论文之后毫无疑问会得到一个副教授(maître de conférences)的职位,因而享有与我相同的社会地位,即“大学教授”。数字最能说明问题:1970年以前开始跟我做工作的十一个学生,一完成工作就得到了副教授职位;而在我的指导下或多或少工作过的那二十来个学生里,没有一个获得这样的职位。诚然,他们之中只有两个人有足够的动力去做国家博士论文(而且两人的论文都很出色)。\par 因此,在这第二个时期,某些暧昧(其深层根源始终被遮蔽着)采取了阶级对抗的形式,采取了对“老板”的不信任(这种不信任被呈现、也被感受为“本能的”“发自肺腑的”)——这就不足为奇了。在那些或多或少当过学生样子的人里,有一个与我的友好关系持续了十来年,没有出现过表面上的对抗插曲,却仍带着同一种暧昧,表现为一种不信任的态度,被“储备”在明显的善意之后。说实话,我从未被这种例行公事式的“不信任”骗过;它在我看来,尤其是这位朋友认为应当给自己找的一个理由,好让自己不去冒险越出他选定的、界限分明的领地——无论在职业生活里,还是在整个生活里都是如此。可这件事他本来可以自由地去做,没有人(顶多只有他自己!)会要他交代。\par 这三个例子,也是我整个教书经验中仅有的、学生(或多少像学生的人)与我之间的关系中某种暧昧以“阶级态度”表现出来的情形。当这种态度出现在大学“团体”内部的同事之间时,显得格外暧昧:在这样一个团体里,双方都享有与常人处境相比极为悬殊的特权,而这些特权使级别(以及薪金)的差别显得相对微不足道。我还注意到,一旦当事人自己也被提升到前一天他还在指责别人的那种处境,这些态度就像中了魔法一样消失了(而且理所当然!)。\par 此外,在我于数学界内部(常常也在外部)目睹的多数——如果不是全部——冲突局面里,我都辨认出一种类似的暧昧。那些“有位置的人”,无论其级别是否合乎他们的期待(合理与否),都享有相当闻所未闻的特权,那是任何别的职业或生涯都无法提供的。那些没有位置的人,则向往同样的安全感和同样的特权(这并不必然妨碍他们对数学本身感兴趣,有时还做出漂亮的工作)。在眼下的世道里,抢位置竞争激烈,没有位置的人常被当作流浪汉对待:我不止一次感到那种默契——在乐于羞辱人的人与被羞辱的人之间;后者咽下去,缩起来。他苦涩与敌意的真正对象,并不是那个行使了权力的人,而恰恰是他自己:他自己缩了起来,并且把对方奉为那种他乐意行使的权力的持有者。乐于羞辱人的人,也是那个报复并补偿(却永远抹不掉……)一场长期承受、早已深埋并被遗忘的羞辱的人。而默许自己被羞辱的人,是他的兄弟和效仿者:他暗中羡慕对方,并在苦涩里埋下那场羞辱,以及羞辱带给他的、关于他自己的那条卑微的消息。} (23''')。

% ---------------- p.177 ----------------
\origpage{68} 在学生与我之间所有这些冲突的显现里,有一个共同点——我教人做数学家这一行快二十五年了——就是强烈的暧昧。无一例外,对抗都是在事后才表现出来的,往往来得很阴险,出现在一段就其本身而言无可怀疑的好感关系里。我甚至可以说,在所有这些情形里,正如在许多别的、并未表现出明确对抗成分的情形里一样,我这个人在他们身上曾产生、并且仍在产生强烈的吸引力。肯定正是这吸引力本身的力量,也滋养了对抗的力量,并保证了它的延续。在对抗采取强烈反感、愤然排斥形式的情形里,肯定也仍然如此;同样,在另一个相反的极端情形里也是如此:在友谊般的敬重这件必不可少的外衣之下,(机会合适时)表达的是一种潇洒而精心调配过的轻蔑的做作……

说实话,这样的暧昧局面并不限于我与某些学生或旧学生的关系。事实上,在我整个成年生活中,从三十岁起(也就是说,从我母亲去世起),它们一直大量存在。在我的情感生活或夫妻生活里是这样,在我与男人的关系里也是这样,更确切地说,主要是与那些明显比我年轻的男人。我终于明白:我身上有某种东西——是天生的还是后天得来的,我说不清——似乎使我倾向于充当父亲的形象。看来我必定有理想的体格和相宜的振动频率,正好造就一个完美的养父!应当说,父亲这个角色我穿起来就像手套一样合手——仿佛它生来就是
%
% ---------------- p.178 ----------------
\origpage{69}\origpage{70} 我的。我不会去数自己有多少次对另一个人进入这样的角色,双方默契得完美无间。大多数时候,这种父—子或父—女的角色分配始终留在未言说之中,甚至留在无意识之中;但也有些时候,它被或多或少清楚地说了出来。在某些情形里,我甚至并没有进入什么游戏,却已经成了父亲的样子——我确信如此——对自己意识层面与无意识层面正在酝酿的事情都一无所知。

我第一次察觉自己扮演养父的角色,是在1972年,“生存与生活”时期,当时我突然面对一位年轻朋友的强烈排斥态度。(有趣的巧合:那是个脱离了常规的数学系学生!)我对第三者的某种举止让他失望了。我想,我本来会毫不困难地承认他的失望是有道理的,承认我在那件事上缺少慷慨——但那份反应的激烈,当时简直把我打懵了。那就像一阵突然腾起的猛烈仇恨之火,而它几乎马上又落了下去,因为已经清楚:他并没有真正把我掀下马。(只差一点点,但这一点我没有让他看出来……)。我不知道当时怎么会直觉到:他把我这个人——被他恰如其分地理想化了的——当成了他与父亲之间未解决冲突的投射对象。这个突然的直觉后来被忘掉了,却没能阻止我在此后许多年里仍然带着同样的确信继续进入父亲的角色,丝毫没有提防。当然,后来当我面对冲突的征兆——阴险的或激烈的——时,也总是同样痛苦的惊讶,不敢相信自己的眼睛,也不敢相信别的。

正是经过一场持续六七个月的、关于我父母生平的紧张独处工作——它让我在一种意想不到的光照下看见他们这个人——我才明白:这种养父的角色是虚幻的;它要替代一个确实存在、却被宣布为(哪怕只是通过默契)“失职”的真正父亲(而且要替代得更好,这是预先就说定了的!)。那是在帮别人回避冲突,把冲突从它所在的地方——比如说,在他与父亲的关系里——转移到第三者(此处即我本人)身上,而第三者与它全然无关。自从那场沉思(它发生在1979年8月到1980年3月)之后,我对自己保持警觉,不再闭着眼睛听任自己那份不合时宜的父性天职发作。这并没有阻止那种虚假局面重演(比如我与那个不得不停止合作的学生之间的关系)——但我相信,如今已没有我这方面的默契。

如果把那个合法期待落了空的学生放在一边,那么对我来说毫无疑问:在我遭遇学生或旧学生对抗的所有其他情形里,那都是同一种“与父亲的冲突”原型的重演:那个既受仰慕又受畏惧、既被爱又被恨的父亲——那个要去对抗、去战胜、去取代、也许还要去羞辱的“他”……但也是那个你暗中想成为的人——夺走他的一股力量,据为己有——另一个自我,被畏惧、被憎恨、被逃避……

\section{(30) 敌父(2)}

在某些旧学生与我之间制造出对抗的,并不是1970年那个大转折——当时的背景是一片牧歌式、无云的过去。它只是让一些对抗变得可见,而那些对抗本来很难在一种典型的、更为循规蹈矩的老板—学生(或前老板—前学生)关系框架里表达出来。我猜想,这类冲突在科学界想必并不罕见,只是它们通常表现得比我所卷入的那些关系更迂回、更难辨认。

回想起来,我终究不觉得在与学生的那些关系里,我曾那么倾向于进入父亲的角色——甚至,我抓不住一个哪怕多少朝着那个方向的记忆。

% ---------------- p.179 ----------------
\origpage{71} 就我这个人而言,我似乎觉得,我投入一段师生关系的几乎全部精力,就是我也投入数学、投入实现一个宏大纲领的那份精力。在第一个时期,我只看到一个例子:我身上对某个学生这个人的兴趣——一种亲近或好感的性质——其力量堪与(如果不说是等同)数学兴趣相比。但即便在那个例子里,我也不觉得我对他进入了父亲的角色。至于我对他、或对其他学生,在某个层次上可能施加的影响,那类事情在我与学生的关系里是我从不留意的。(即便今天,我也倾向于不留心,无论对近几年跟我做过工作的学生,还是对别的人。)当然,在所有这些情形里,学生与我之间的关系绝不是“对称”的:至少在师生关系持续的那段时间里(而且多半还要更久),一个学生在我生命里的分量,与我必定在他生命里占据的分量无法相比;关系在我身上和他身上所调动的心理力量也无法相比。除了那五六个以清晰可辨的对抗征兆表现出这些力量的情形,我意识到:在二十多年教书生涯里,我那些不同的学生、后来是旧学生与我之间关系的性质,对我来说始终是个完全的谜!再说,探究这些谜也不是我分内的事,那更多是各人自己那一份的事。不过,既然关心起自己这个人来,也许还有些比“他与旧老板关系的来龙去脉”更灼人的东西可以看。无论如何,尽管我在学生面前没有表现出任何进入父亲角色的倾向,我多半还是常常在他们眼里多少充当了养父——鉴于我前面谈过的自己那份特殊的心理“剖面”,也鉴于一种局面的内在动力:在那种局面里,我起码不能不显出一个长辈的样子。

无论如何,在我提到的几个例子里,学生与我之间关系的这种特殊色调,在我看来毫无疑问。在职业生活之外,还有许多别的例子:不管有没有我的默契,我显然在比我年轻的男人或女人面前充当了养父——他们被这个人吸引,首先与我以相互的好感相连,而绝无亲缘关系。至于我自己的孩子们,我对他们的父性之弦是强烈的;从很小的时候起,他们在我生命里就占据着重要的位置。奇怪的讽刺是:我的五个孩子,没有一个接受“有我做父亲”这件事。在我能够近距离了解的四个孩子的生活里——尤其是近几年——他们与我关系中的这道分裂,是他们自身内在一道深刻分裂的反映;尤其是一种拒绝:拒绝他们身上一切与我这个父亲相像的东西……但这里不是探究这道分裂之根的地方:那些根既扎在一段被撕裂的童年里,也扎在我的童年、我父母的童年里;同样也扎在母亲的童年、她父母的童年里。这里也不是衡量其后果的地方——在他们自己的生活里,或在他们孩子的生命里……

\section{(31) 使人泄气的权力}

为了结束这趟粗略的巡游——穿过1948到1970年间我在数学界里的种种关系——我还得谈谈我与更年轻的数学家的关系:他们或多或少是新手,因而严格说来没有“同事”身份,而我也不曾对他们扮演“老板”的角色。所以,这里说的是些年轻研究者:我在IHÉS的讨论班上与他们相遇一两年,或在哈佛及别处的这类课程、讨论班上,有时也通过
%
% ---------------- p.180 ----------------
\origpage{72}\origpage{73} 通信——比如我收到一位年轻作者寄来的工作,他期待我给出评论,肯定也期待鼓励。

与初入门研究者的关系,属于一个不如“学生们的老板”那么显眼的角色,但同样重要——这是我后来才明白的。那个年代,我并不像这六七年以来那样意识到:对一位有名望的数学家来说,这个角色代表着一种相当大的权力。首先是鼓励、激励的权力:它既存在于明显出色的工作那里(尽管那工作也许因表述笨拙或“行内功夫”不足而吃亏),也存在于仅仅扎实的工作那里;它甚至存在于那种只代表一份极微薄贡献的工作那里——按一位正当壮盛、手段强劲、对主题经验老到、见闻广博的前辈的标准,那贡献甚或微不足道,乃至等于零。只要送到我们手里的工作是认真写成的——这通常从头几页就看得出来——鼓励的权力就在那里。

而使人泄气的权力同样存在,而且无论面对什么样的工作,都可以随意行使。这就是柯西对伽罗瓦、高斯对雅可比行使过的那种权力——它可不是从昨天才有的,杰出而令人畏惧的人物一直都在使用它!历史之所以把这两桩事传给我们,是因为付出了代价的这两个人有足够的信念和自信继续走自己的路,尽管那些当时在数学界呼风唤雨的人怀着的权威毫无善意。雅可比找到了一份杂志发表他的想法,伽罗瓦找到的则是他最后一封信的纸页,权当“杂志”。

今天,对一个无名或不太有名的数学家来说,要让人知道自己,肯定比上个世纪更难。而有名望的数学家的权力,不只在心理层面,也在实践层面。他有权接受或拒绝一篇工作,也就是说:给不给它发表的支持。对也好错也好,我觉得“在我那个时代”,五六十年代,拒绝并非不可上诉——如果工作呈现出“值得注意”的结果,它还有机会得到另一位权威的支持。今天肯定不再如此:如今甚至很难找到一位有影响的数学家,肯(以他乐意采取的心情)去通读一篇属于他专业领域的工作——如果作者尚未获得名声,或者没有一位他认识的同行推荐。

近几年,我曾见有影响、出色的数学家行使他们使人泄气的权力,拒绝这个、拒绝那个:既针对某些显然应当做出来的扎实工作,也针对某些规模可观、清楚显示出作者力量与原创性的工作。有好几次,这样行使自己任意权力的人,竟是我的一个旧学生。这无疑是我作为数学家的生涯里经历过的最苦涩的经验。

但我又跑题了!我本来要考察的是:在我诚心诚意地扮演“知名数学家”角色的那些年月里,我是怎样使用自己手中那份鼓励人和使人泄气的权力的。我还该补充:在1970年以后我的科学活动所处的那个更卑微的层次上——在外省大学里作为普通教师中的一员——这种权力并没有因此就不存在,无论是对我的学生,还是(说实话,很少)对偶尔通信的人。但就我眼前的题目而言,只有我数学生涯的第一个时期才是要紧的。

至于我与学生的关系,从我带过的第一个直到今天,我相信可以毫无保留地说:我做了力所能及的一切去鼓励他们
%
% ---------------- p.181 ----------------
\origpage{74} 在自己选定的工作上往前走\fnoteP{12}{(23$^{iv}$)}{\emph{一次教学的失败(1)}\par 这些字句写下之后,我有机会与1970年之后的两位旧学生谈过,想同他们一起探究我在蒙彼利埃大学研究层次上教学失败的原因。他们对我说:我那种倾向——低估同化那些对我熟悉、对他们却不熟悉的技术对他们可能意味着的困难——在他们身上产生了令人泄气的效果,因为他们感到自己始终达不到我对他们的期望。此外(这件事在我看来意义还要大),他们有时感到受了挫伤:当我把袖子里藏着的一个成形陈述“卖给他们”、向他们泄底的时候,而不是留给他们乐趣,让他们在已经十分接近的时刻靠自己的办法把它发现出来。那之后,他们剩下的就只有做“练习”了(而那并不特别让他们着迷)——也就是证明那个陈述。我先前在一则注(注21)里指出过的自己身上那份“缺乏慷慨”,正落在这一处,当时没有就此多说。主要正是这类扫兴之事,构成了我在他们两人身上——在那么出色的开端之后——研究兴趣消失这件事上的个人贡献。\par 我意识到,我在1970年以前并不比之后更慷慨。如果说那时我没有遇到同样的困难,那无疑是因为:那个时期来找我的那类学生,有足够的动力,甚至能从一场“长长的练习”里品出魅力——那是学习本行、并一路学到一大堆东西的机会;同样,对于一个我“泄了底”的起步陈述,他们也能靠自己的办法从中抽出一大群远远超出第一个的陈述。当我换了教学地点,我在向新学生提议的思考主题上做了必要的调整,办法是选择那些可以被直接直觉把握的数学对象,不依赖任何技术行囊。但这必要的调整本身并不够,因为禀性上的差别(我的新学生与从前的学生相比)比单纯的行囊差别还要重要。这还印证了前面(第25段开头)说过的一点:我在“导师”这个角色上有某种不足;这不足在我教学生涯的第二个时期,比在第一个时期显露得强烈得多。} (23'''')。即便在今天,老板与学生的关系里恐怕也很少是别的样子,尤其当这位老板拥有手段,能够培养出色的学生,并借他们的协助开垦大片待耕的荒地。几乎难以置信、然而千真万确的是:甚至存在这样一种极端情形——一位声望卓著的老板,乐于把曾经激励过他自己年轻时代的数学激情,在才华出众的学生身上熄灭掉。

但我又跑题了!现在要考察的,是我与那些不是学生的年轻研究者的关系。在那样一些关系里,名人身上的自我中心力量,不那么容易把他推向鼓励的方向,因为那个找上他的无名年轻人的成功,对他的荣耀毫无增益,或增益甚微。恰恰相反,我认为:在缺乏真正善意的情况下,仅凭自我中心力量的运作,几乎无一例外地会把人推向相反的方向——去行使使人泄气、拒绝的权力。在我看来,这恰恰就是那条在社会一切部门都能看到的普遍法则:证明自己重要的自我中心欲望,以及满足它时伴随的秘密快乐,通常在这样一种情形下更强、更受享用——手里的权力找到机会去造成别人的失望,甚至羞辱,而不是相反。这条法则在某些特殊处境里表现得格外粗暴,比如战争,比如集中营的世界,监狱或精神病院的世界,或者干脆就是我们这样一个国家里任何人都可能进的医院……但即便在最日常的处境里,我们每个人都遇到过印证这条法则的态度与行为。对这些态度的矫正,首先来自文化的矫正:在某个圈子里,关于什么算“正常”或“可接受”的行为有一种共识;其次则是非自我中心性质的力量,比如对某个特定的人的好感,或者有时,一种自发的善意,甚至与它所指向的人无关。这样的善意无疑是稀罕之物,无论你去哪个圈子里找。至于数学圈里的文化矫正,在我看来,在过去的二十年里已经被大大侵蚀了。至少在我所熟悉的
% ---------------- p.182 ----------------
\origpage{75}\origpage{76} 那些圈子里,情况肯定如此。

我执意跑题,真是没办法!我本来要说的不是一篇关于这个世纪的议论,而是一场对我自己、对我与那些或多或少刚刚入门、又不是我学生的研究者之间关系的沉思。我不认为我提到的那条“法则”在那些关系里找到过表现的机会。由于一些不必在此考察的原因,看来自我中心的力量,虽然在我身上和在任何人身上一样强,在我一生中却没有走那条路去损害别人(除了几桩追溯到童年的情形)。我甚至敢说——既然有机会考察这件事——我面对他人的基本调性是一种善意的调性:能帮的时候就帮,能缓解的时候就缓解,有能力鼓励的时候就鼓励。即便在与那位“不知疲倦的朋友”如此深刻分裂的关系里,我身上的自负也从未把我误导到这样的地步:让我想到(哪怕是无意地)去伤害他。(我本来是有可能那样做的,而且当然会“怀着世上最安宁的良心”。)我相信,在大多数情形里,这种普遍的善意(哪怕只是浮在皮肤表面的一点点)也标记了我在数学界里的关系,包括与那些虽不在学生之列、却可能需要我支持或鼓励的初入门数学家。

我相信,至少在五十年代、直到六十年代初,无一例外都是如此。我觉得至少在那个时期,这份善意并不限于像广中平祐(Heisuke Hironaka)或迈克·阿廷(Mike Artin)那样明显出色的年轻人(当时还没有任何名声证明他们的本领)。但它在六十年代也许或多或少地淡去了,在自我中心力量的作用下。关于这一点,若有任何见证传到我这里,我会格外感激。

我的记忆只交给我一个确切的例子,我这就说;而在这个例子之外,就是那团著名的“雾”——它不凝结成任何别的确切事例或事实,却交给我某种内心的态度。当另一位数学家“踩到我的地盘上”、却装作不问我一声,好像他在自己家里似的——这个毛头小子!——我会感到某种恼火。那多半是年轻人的情形:不大在行,却碰巧重新找到一些我多年来、而且是从高处就知道的东西,有时还偏偏在一些很特别的情形里。我想,这种事不至于常常发生,但也许还是有两三次,也许是四次,我说不准。如我刚才所说,我只记得一个具体的例子,也许是因为那局面与同一位年轻数学家反复出现过好几次,形式各不相同。我可以说:这位年轻研究者——他所属的大学在国外——在各方面都极其得体:他把自己刚做好的工作寄给我,而我理应是消息最灵通的人。每一次,我都反应得极为冷淡,理由就是我说的那个。我甚至说不准自己有没有坦率地告诉他:他做的东西我早就知道了,因此他不至少在前言里向我行个小小的屈膝礼就发表,让我很不痛快。当然,假如他是我的学生,这种作者的虚荣就不会起那么大作用:一方面因为与学生之间已经建立起好感,另一方面因为无论如何不言而喻,学生的工作里也有老板的想法,除非另有说明!我相信这局面与同一位研究者发生过两次,也许甚至三次;每一次我的态度都一样冷淡,一样令人泄气。如果我记得不错,我从未答应过推荐这位研究者的工作到某家杂志发表,也从未答应参加他的论文答辩委员会(我记得这个问题曾被提起过)。几乎可以说,我像是打定主意挑他当出气筒。最妙的是,
%
% ---------------- p.183 ----------------
\origpage{77} 他的工作每一次都完全站得住脚——我相信写得用心,我也没有任何理由假定那些想法不是他自己找到的:那些想法当时还不怎么满街乱跑,只是(或多或少)“众所周知”于一小撮在行的人之间,比如塞尔、卡蒂耶、我,还有一两个人。我不明白的是:这位年轻同事(他当然最终有了论文,也有了一个受之无愧的职位)怎么就没厌烦一次次找上我——那个每次“给他冷脸”的人,而且他显然从未为此怨过我。我还记得有一次,他对我的迟疑表示惊讶;显然他不明白发生了什么。要是他等着我的解释,那可有他难受的!他有一张好看的脸,有点像古典希腊的,非常年轻——线条相当柔和、安详,透出一种内在的平静……现在,我第一次试着去捕捉他这个人和他的相貌所散发的印象,忽然意识到:他真的很像我说过的那位“不知疲倦的朋友”;他们本可以是兄弟——这位与我同龄的朋友,调子是含笑的那种;而这位研究者,比他年轻二十岁,调子略略低沉,却绝不忧伤。这种相似未必没有起过作用:我把一份在另一个人身上找不到机会表达的轻蔑,投射到了这一位身上——因为另一个人以那样忠实的友谊征兆把我解除了武装!而我也确实得长出一副多么厚的甲壳,才没有被这个必定招人喜爱的年轻人的明显诚意和把事情做好的意愿所打动:他不厌其烦地再来,而我连一个微笑都不肯赏他!

\section{(32) 数学家的伦理}

我昨天讲的那个例子,现在我终于肯费心把它白纸黑字记下来,它在我看来具有相当大的意义,在某些方面比我先前讲到的另外三个例子(无疑同样典型)更大——在那三例里,自负的力量深深地扰乱了我身上一种自然的善意与尊重的态度。这一次,我动用了一个实实在在的权力位置(同时像所有人一样,装作不知道这份权力),用它去使一位善意的研究者泄气,拒绝一篇本应发表的工作。这就是人们所说的滥用权力。它并不因为不落入刑法某一条的制裁范围就不那么昭然。幸而当时的形势不如今天严酷,所以这位研究者能够——我相信,没费太大力气——在几位比我更善意的同行的支持下把工作发表出来;他作为数学家的生涯,也没有因为我那滥用权力的行为而受到严重扰乱,更不用说被毁掉。事后我为此高兴,却并不想把它当作“减轻情节”。在更严酷的形势里,我也许会更当心——但这只是一个推测,在这里没有多少关系。不过我仍相信自己可以说:我心里并没有暗藏的恶意,没有由我说的那种恼火引起的害人之心。我对那恼火的反应是“发自肺腑的”,对自己没有丝毫批评的意愿,更没有丝毫意愿去看一眼自己心里发生了什么,哪怕只是看一眼我的反应可能在另一个人的生命里有什么分量。我没有衡量自己手中的权力;与这权力相随的责任这个念头,在那段关系里从未掠过我的心头。这是一桩不负责任行为的典型,就像街头巷尾随处可遇的那样,在科学界和在别处一样。

我记住的这个同类情形中唯一的例子,也许只是几个类似情形中的一个极端例子。引发一种毫无善意的态度的,是虚荣心受到冒犯后的恼火:它
%
% ---------------- p.184 ----------------
\origpage{78}\origpage{79} 容不得“随便哪个初来乍到的人”僭取权利,闯进禁猎区,打走一点只属于这些地盘之主的小猎物……这种恼火有的是现成的合理化,而且派头更高贵,这可想而知。被摆出来的不是我这卑微的个人——不!——而是对艺术、对数学的爱;这个年轻人连“是个天才”这条借口都没有,倒更像笨手笨脚的那一类,他会把一切都弄糟,我们可倒霉了;他要是把事情做得比我好倒也罢了,可我筹划的那些漂亮安排全泡汤了,说真的,总得有点不知羞才行……!而始终如水印般透出的,是那支功绩主义的主导动机:只有最好的那些人(比如我)在我的地盘上才有居留权,或者那些把自己置于这类人中某一位保护之下的人!(至于那种不那么常见的情形——确实是另一位大人物踩到我的地盘上——那又是另一回事了:一天自有一天的烦恼!)在这个具体例子里,还有(对此我如今已不大有疑问)另一股朝同一方向起作用的力量,而它完全是无意识的;这股力量在我与起步时那位不知疲倦的朋友的关系里,已经起过很强的作用:对某一类人的排斥自动机制——那类人不合我从母亲那里接过来的“男子气”标准。但这个情形,虽然对理解我自己有它的意义和趣味,与我眼前的题目却关系不大:我的题目是要在我自己身上,在我还属于某个圈子时曾有过的态度和行为里,找出我今天在那个圈子里所看到的深刻败坏的典型征兆。

我刚考察的这个例子,之所以在我看来比我其他那些缺乏善意与尊重的例子意义更大,是因为正是在它这里,数学家这一行某种基本的伦理遭到了侵犯\fnoteP{13}{(24)}{我想说的这种伦理,同样适用于任何别的围绕研究活动形成的圈子;在那里,把自己的结果公之于众、并从中获得承认的可能性,对每个成员的社会地位、甚至对作为该圈子成员的“存活”,都是“生死攸关”的问题,连带对他和他的家庭意味着一切后果。} (24)。在我起步时接纳我的那个圈子里——也就是布尔巴基及其亲近者的圈子——我想说的这种伦理大体上始终是隐含的,但它毕竟在场、有生命,而且(在我看来)是不可触动的共识的对象。据我记忆所及,唯一把它用清楚明确的话说出来的人是迪厄多内,那无疑是我在南锡做客的最初几次之一。他也许在别的场合又回到这一点。显然他觉得这是一件重要的事;而我当时必定感受到了他对此的看重,因为三十五年后的今天我还能记得。仅仅凭借我那些年长的前辈这个群体的道德权威,以及显然表达了群体共识的迪厄多内,我必定已经默默地把这种伦理认作自己的,却从未给它一刻思考,也从未明白它的重要性何在。说实话,我甚至不会想到:对它加以思考也许有用——因为我早已确信,我的父母和我自己这个人,各自都是伦理态度的完美化身(或者相差无几),负责、全面,而且经得起一切考验\fnoteP{14}{(25)}{\emph{职业伦理的共识——兼及信息控制}\par 除了与迪厄多内的那次谈话,在我作为数学家的生涯里,我记不起有哪次谈话——我参与的或旁听的——谈到这一行的伦理、谈到行业成员之间关系中的“游戏规则”。(这里我排除那些关于科学家与军事机构合作的讨论,它们发生在七十年代初“生存与生活”运动前后。那些讨论其实并不涉及数学家之间的关系。我在“生存与生活”里的许多朋友,包括谢瓦莱和盖吉,都觉得:我当时——尤其起初——把重点放在这个我格外敏感的问题上,使我远离了更本质的日常现实,正是我在眼下这份反思里考察的那类现实。)在学生与我之间,从未谈过这些事。那种心照不宣的共识,我想只限于这一条规则:不把别人的、自己可能已经知道的念头说成是自己的。这共识,在我看来自古就存在,而且至今没有任何科学圈子里的人质疑过。但若缺少另一条补充规则——保障每个研究者有可能让人知道自己的想法和结果——第一条规则就只是一纸空文。在今天科学界,处于声望与权力地位的人,对科学信息握有一种任意处置的控制权。在我所熟悉的那个圈子里,这种控制不再被迪厄多内所说的那种共识所缓和;而那共识,也许从未存在于他充当代言人的那个狭小团体之外。处于权力地位的科学家,几乎能收到他认为值得收到的一切信息(往往还超出),而且对其中很大一部分,他有权阻止发表,同时坐享其利,把那些信息斥为“没有意思”、“多少已知”、“trivial”等等……我在注(27)里还要回到这个局面。} (25)。

% ---------------- p.185 ----------------
\origpage{80} 再说,迪厄多内并没有对我长篇大论——那既不是他的风格,也不是他布尔巴基的任何朋友们的风格。他多半是顺口提起,像提起一件理应不言自明的事。他只是强调一条再简单不过、表面上全然无害的规则,就是这样:任何人发现了一个值得注意的结果,就应当有权利、也有可能把它发表,唯一的条件是这结果尚未成为某份出版物的对象。所以,即便这个结果为一人或数人所知,只要他们不曾费心把它白纸黑字写下来并发表,从而使它可供(哼!)“数学共同体”使用,那么任何别的人(言下之意:包括那位著名的“初来乍到者”)用自己的办法找到这个结果(言下之意:不管他用的是什么办法、什么观点和光照,也不管这些在那些理应比他更在行的人看来是否“褊狭”……),就应当有可能按自己的办法和光照把它发表。我记得迪厄多内还补充说:这条规则若不被遵守,就会给最恶劣的滥用敞开大门——很可能正是在那个场合、经他的口,我才知道了高斯拒绝雅可比工作那桩历史案例——借口是雅可比的想法他早已知道。

这条如此简单的规则,是对迪厄多内身上(以及布尔巴基其他成员身上)那种“功绩主义”态度——正如在我身上一样——的根本矫正。遵守这条规则,是一种廉洁的保证。我很高兴能够说,从至今传到我这里的一切来看,这种根本的廉洁在布尔巴基最初群体的每一位成员身上都保持完好\fnoteP{15}{(26)}{布尔巴基的“创始成员”是亨利·嘉当、克劳德·谢瓦莱、让·德尔萨特(Jean Delsarte)、让·迪厄多内、安德烈·韦伊。他们全都健在,只有德尔萨特例外——他在五十年代英年早逝,而那个时期,这一行的伦理总体上还受到尊重。} (26)。我注意到,对于曾经属于布尔巴基群体或圈子的别的数学家,情况并不都是这样。而在我自己这个人身上,它并没有保持完好。\par{}重读本文时,我曾想删掉这一段:在里头我可能给人这样的印象,像在颁发“廉洁”(或不廉洁)证书,而当事人对此毫不在乎,也不是我该做的事。这段文字可能引起的保留,肯定是有道理的。但我还是把它留着,为的是见证的真实性,也因为这段文字确实再现了我的感受,即便这些感受并不合宜。

迪厄多内用再平实不过的话向我讲的那种伦理,作为一种特定圈子的伦理,已经死了。或者说,那个圈子本身也死了,与那构成它灵魂的廉洁一同死去。这份廉洁在某些孤立的人身上保存了下来;在另一些它曾经败坏的人身上,它重现了,或者将会重现。它在某个人身上出现或消失,是各人精神历险的关键插曲之一。但这历险得以展开的舞台,已经深刻地改变了。一个曾经接纳我、被我认作自己的、我曾暗自引以为傲的圈子,不复存在了。使它有价值的东西,在我自己身上死去了,或者至少被另一种性质的力量侵入并取代——而那远在规范它的默会伦理于使用惯例与信仰表白中被公开否弃之前。如果我此后能够惊讶、能够愤慨,那是出于故意的无知。从这个曾属于我的圈子里回到我身上的东西,本来有一个关于我自己的消息要带给我,而我乐意把它一直回避到今天。

% ---------------- p.186 ----------------
% (本页仅存书眉“8. Maître et élèves”,无正文)

